% ================================================================
%  PDL --- D34
%  Born Rule from (A)∧(B)-Admissible Amplitudes in the PDL Framework
%  C. Laubscher, 2026
% ================================================================
\documentclass[11pt,a4paper]{article}
\usepackage{mathtools}
\usepackage[utf8]{inputenc}
\usepackage[T1]{fontenc}
\usepackage[british]{babel}
\usepackage{amsmath,amssymb,amsthm}
\usepackage{mathrsfs}
\usepackage{graphicx}
\usepackage[colorlinks=true,citecolor=blue,linkcolor=blue,urlcolor=blue]{hyperref}
\usepackage[margin=2.5cm]{geometry}
\usepackage{booktabs}
\usepackage{enumitem}
\setlist[itemize,1]{label=---}
\usepackage{csquotes}
\usepackage{xcolor}
\usepackage{mdframed}

% ── Theorem environments ──────────────────────────────────────────────────────
\newtheorem{theorem}{Theorem}[section]
\newtheorem{proposition}[theorem]{Proposition}
\newtheorem{lemma}[theorem]{Lemma}
\newtheorem{corollary}[theorem]{Corollary}
\theoremstyle{definition}
\newtheorem{definition}[theorem]{Definition}
\newtheorem{remark}[theorem]{Remark}
\newtheorem{axiom}[theorem]{Axiom}

% ── Macros ────────────────────────────────────────────────────────────────────
\newcommand{\PDL}{\textnormal{PDL}}
\newcommand{\Rsurf}{R_{\mathrm{surf}}}
\newcommand{\Rtot}{R_{\mathrm{tot}}}
\newcommand{\rval}{r_{\mathrm{val}}}
\newcommand{\Rsea}{R_{\mathrm{sea}}}
\newcommand{\hb}{\hbar}
\newcommand{\Ree}{R_e}
\newcommand{\vp}{\varphi}
\newcommand{\aPDL}{\alpha_{\PDL}}
\newcommand{\Kfour}{K_4}
\newcommand{\Hspin}{\mathcal{H}_{\mathrm{spin}}}
\newcommand{\emix}{\varepsilon_{\mathrm{mix}}}
\newcommand{\ecoh}{\varepsilon_{\mathrm{coh}}}
\newcommand{\Tplus}{\mathcal{T}_+}
\newcommand{\Tminus}{\mathcal{T}_-}
\newcommand{\pket}[1]{|#1\rangle}
\newcommand{\pbra}[1]{\langle#1|}

% ================================================================
\begin{document}

\title{Born Rule from $(A)\wedge(B)$-Admissible Amplitudes in the \PDL\ Framework}
\author{C.~Laubscher\thanks{Independent researcher, Switzerland. \texttt{cedric.laubscher@gmail.com}}}
\date{\small \today}
\maketitle

% ================================================================
\begin{abstract}
The Projective Dynamic Logo (\PDL) framework derives physical constants and dynamical laws from four axioms on finite signed graphs. Paper~D32 of the corpus introduced complex amplitudes $\psi_n$ for $(A)\wedge(B)$-admissible transitions of $\Kfour$, but the probability interpretation $|\psi_n|^2$ was not derived from the axioms. The present paper establishes the connection between the \PDL\ coherence structure and Born's rule at two distinct epistemic levels, with honest documentation of what is proved and what remains open.

\emph{Level~1 (proved).} The coherence violation cost $\ecoh(\sigma;\theta,\psi) = |\pbra{-\theta}\pket{\psi}|^2$ is constructed from the integer-valued indicator $\chi(\tau;\sigma) \in \{+1,-1\}$, which counts whether a mixed triangle is coherent or violated under the sign configuration $\sigma$ --- a purely combinatorial quantity. The identity $P_+ = 1 - \ecoh = |\pbra{+\theta}\pket{\psi}|^2$ follows algebraically from normalisation. Five Gleason axioms (normalisation, global phase invariance, rotational covariance, branch symmetry, repeatability) are stated in full and verified for this assignment to $\leq 1.1\times10^{-15}$ over 20\,000 test cases. A Gleason-type uniqueness theorem is stated (self-contained in Section~\ref{sec:gleason}) and applied: $P_+ = |\pbra{+\theta}\pket{\psi}|^2$ is the \emph{unique} admissible probability assignment. This closes open problem~OP5 of D32 at Level~1.

\emph{Level~2 (open problem~OP1).} The identification $\ecoh = |\pbra{-\theta}\pket{\psi}|^2$ requires that the triangle weights $\alpha_\tau(\theta,\psi)$ be distributed non-uniformly according to the inner-product geometry of $\Hspin$. A diagnostic result (Proposition~\ref{prop:heaviside}) shows that uniform weights lead to a Heaviside step function, not Born's rule, pinpointing exactly what remains to be derived from the \PDL\ axioms alone.

The document is self-contained: the key results from the related \PDL\ corpus documents~\cite{Born_PDL_2026,Gleason_PDL_2026} on which D34 relies are stated explicitly in Sections~\ref{sec:born_scheme} and~\ref{sec:gleason}, so that a reader without access to those documents can verify the argument independently.
\end{abstract}

% ================================================================
\section{Introduction}
\label{sec:intro}

The \PDL\ programme reconstructs fundamental physics from four axioms on finite signed graphs~\cite{PDL_Main_2026,PDL_Intro_2026}. The proton is characterised by the unique integer quintuplet $(n_u, n_d, \rval, \Rsea, \Rtot) = (24, 28, 930, 10087, 11017)$, and the minimal electron prototype is the complete signed graph $\Kfour$ on four vertices and six edges, with relational budget $\Ree = 6$~\cite{Combinatorial_Proton_v2_2026}.

Paper~D32 derived the non-relativistic Schr\"odinger equation from the $(A)\wedge(B)$ stability criterion of D29 and the Compton pulsation of $\Kfour$~\cite{PDL_D32_2026}. Complex amplitudes $\psi_n$ are introduced as a coarse-grained description of the $(A)\wedge(B)$-admissible dynamics; the transition coefficient $b = -i\hb\Delta t/(2m_e(\Delta x)^2)$ is derived from condition~(B) without free parameter. However, the probability interpretation $P_n = |\psi_n|^2$ --- Born's rule --- is listed as open problem~OP5 of D32.

Two documents of the \PDL\ corpus address this directly. Document~\cite{Born_PDL_2026} constructs a mixed-triangle measurement scheme in which Born probabilities emerge as basin-ratio weights in the limit of sharp coherence optimisation. Document~\cite{Gleason_PDL_2026} establishes a Gleason-type uniqueness theorem for spin-$\tfrac{1}{2}$: under five operational axioms, Born's rule is the \emph{unique} admissible probability assignment on $\Hspin$. The key statements from both documents are reproduced in full in Sections~\ref{sec:born_scheme} and~\ref{sec:gleason} of the present paper, making D34 self-contained.

The present paper builds the bridge between these results and the $(A)\wedge(B)$-admissible amplitudes of D32. The key object is the coherence violation indicator $\chi(\tau;\sigma) \in \{+1,-1\}$: a deterministic integer function of the sign configuration $\sigma$ of $\Kfour$ and the apparatus. From $\chi$, one constructs the coherence cost $\ecoh$ that equals $|\pbra{-\theta}\pket{\psi}|^2$ without assuming Born's rule. The Gleason uniqueness theorem then closes the argument.

% ================================================================
\section{\PDL\ prerequisites}
\label{sec:prereq}

\subsection{The \texorpdfstring{$K_4$}{K4} closure, pulsation, and condition~(B)}

$\Kfour$ is the complete signed graph on four vertices $\{0,1,2,3\}$ and six edges. A sign assignment $s^{(1)}$ satisfying triangular coherence ($s^{(1)}_{ij}s^{(1)}_{jk}s^{(1)}_{ik} = +1$ for every triangle) is unique up to the global flip. The binary pulsation alternates between $s^{(1)}$ and $s^{(2)} = -s^{(1)}$~\cite{PDL_Main_2026}. There are exactly 8 coherent sign configurations of $\Kfour$, forming 4 pulsation pairs $\{s,-s\}$.

A mixed triangle $(e_i, e_j, p_k)$ with $e_i, e_j \in V(\Kfour)$ and $p_k$ a proton vertex is stable across both half-cycles if and only if~\cite{PDL_D29_2026}:
\begin{align}
\mathrm{(A)}&\quad P_1 = s^{(1)}(e_i,e_j), \label{eq:condA}\\
\mathrm{(B)}&\quad P_2 = -P_1. \label{eq:condB}
\end{align}
This was verified exhaustively over 768 cases with zero counter-example; $P_2/P_1 = -1$ in all 192 stable cases (Table~\ref{tab:condB}, self-contained verification).

\begin{table}[ht]
\centering
\caption{Self-contained verification of the stability criterion (A)$\wedge$(B) for mixed triangles in the \PDL\ framework. The sign $s^{(1)}(e_i,e_j)$ is fixed at $+1$ without loss of generality. Counts over 768 cases (8 coherent $K_4$ configurations $\times$ 6 edges $\times$ $2^4$ cross-sign combinations); 192 stable; 0 counter-examples to the equivalence.}

\bigskip

\label{tab:condB}
\begin{tabular}{cccccr}
\toprule
$P_1$ & $P_2$ & (A): $P_1 = s^{(1)}$ & (B): $P_2 = -P_1$ & $(A)\wedge(B)$ & Stable \\
\midrule
$+1$ & $+1$ & $\checkmark$ & $\times$ & $\times$ & No \\
$+1$ & $-1$ & $\checkmark$ & $\checkmark$ & $\checkmark$ & Yes \\
$-1$ & $+1$ & $\times$ & $\checkmark$ & $\times$ & No \\
$-1$ & $-1$ & $\times$ & $\times$ & $\times$ & No \\
\bottomrule
\multicolumn{6}{l}{\small $P_2/P_1 = -1$ in all 192 stable cases. $(A)\wedge(B) \Leftrightarrow$ stable: 768/768.}
\end{tabular}
\end{table}

\subsection{The spin doublet of the \texorpdfstring{$(4,6)$}{(4,6)} block}

The $(4,6)$ block admits two logically stationary motifs, identified with the spin-up and spin-down basis states $\{\pket{+}_z, \pket{-}_z\}$ of $\Hspin \cong \mathbb{C}^2$. A general state is $\pket{\psi} = a\pket{+}_z + b\pket{-}_z$ with $|a|^2 + |b|^2 = 1$. The measurement basis along direction $\theta$ is:
\begin{equation}
\pket{+}_\theta = \cos(\theta/2)\pket{+}_z + \sin(\theta/2)\pket{-}_z, \qquad \pket{-}_\theta = \sin(\theta/2)\pket{+}_z - \cos(\theta/2)\pket{-}_z.
\label{eq:basis_theta}
\end{equation}
The inner products are $\pbra{+\theta}\pket{\psi} = a\cos(\theta/2) + b\sin(\theta/2)$ and $\pbra{-\theta}\pket{\psi} = a\sin(\theta/2) - b\cos(\theta/2)$.

% ================================================================
\section{The mixed-triangle Born scheme (self-contained)}
\label{sec:born_scheme}

The following summarises the key construction of~\cite{Born_PDL_2026}, reproduced in full here so that the reader can verify D34 without access to that document.

The apparatus is coupled to the $(4,6)$ system via $N_T$ mixed triangles on the active surface $\Rsurf = 310\vp$~\cite{Effective_Sigma_PDL_2026}. These triangles are partitioned into two families by their apparatus sign:
\begin{equation}
\Tplus\,\text{(apparatus sign $+1$)}, \qquad \Tminus\,\text{(apparatus sign $-1$)}.
\label{eq:partition}
\end{equation}
The geometry of the partition depends on the measurement direction $\theta$; the active-surface fractions are:
\begin{equation}
\frac{N_T^+}{N_T} = \cos^2(\theta/2), \qquad \frac{N_T^-}{N_T} = \sin^2(\theta/2).
\label{eq:fractions}
\end{equation}
A global sign configuration $\sigma$ on $\Kfour \cup G_{\mathrm{app}}$ is weighted by $w(\sigma;\theta,\psi) = \exp[-\lambda\,\emix(\sigma;\theta,\psi)]$ with $\lambda \gg 1$, where $\emix$ is the mixed-triangle coherence cost (Definition~\ref{def:emix} below). The two outcome basins are:
\begin{equation}
A = \{\sigma : \text{predominantly $\Tplus$-coherent}\}, \qquad B = \{\sigma : \text{predominantly $\Tminus$-coherent}\}.
\label{eq:basins}
\end{equation}
The outcome probability is $P_+(\theta;\psi) = W_A/(W_A + W_B)$ where $W_A = \sum_{\sigma \in A} w(\sigma;\theta,\psi)$ and $W_B = \sum_{\sigma \in B} w(\sigma;\theta,\psi)$. In the limit $\lambda \to \infty$, under the mild apparatus symmetry conditions of~\cite{Born_PDL_2026} (exchange symmetry between $\Tplus$ and $\Tminus$ and SU(2)-covariance of the apparatus interface), the weights satisfy:
\begin{equation}
\frac{W_A}{W_B} \approx \frac{|\pbra{+\theta}\pket{\psi}|^2}{|\pbra{-\theta}\pket{\psi}|^2},
\label{eq:weight_ratio}
\end{equation}
yielding $P_+(\theta;\psi) = |\pbra{+\theta}\pket{\psi}|^2 = |a\cos(\theta/2) + b\sin(\theta/2)|^2$. Table~\ref{tab:born_check} gives a self-contained numerical verification of equation~\eqref{eq:weight_ratio} for four representative states.

\begin{table}[ht]
\centering
\caption{Numerical verification of Born's rule from the coherence complement: $P_+ = 1 - |\pbra{-\theta}\pket{\psi}|^2 = |\pbra{+\theta}\pket{\psi}|^2$, for four representative states over 200 measurement angles $\theta \in [0,\pi]$. All residuals are at floating-point precision ($\leq 5.55\times10^{-16}$), confirming the algebraic identity.}

\bigskip

\label{tab:born_check}
\begin{tabular}{lr}
\toprule
State $\pket{\psi}$ & $\max_\theta |P_+ - |\pbra{+\theta}\pket{\psi}|^2|$ \\
\midrule
$\cos(\pi/8)\pket{+}_z + \sin(\pi/8)\pket{-}_z$ & $3.33\times10^{-16}$ \\
$\tfrac{1}{\sqrt{2}}(\pket{+}_z + \pket{-}_z)$ & $4.44\times10^{-16}$ \\
$\tfrac{1}{\sqrt{2}}(\pket{+}_z + i\pket{-}_z)$ & $4.44\times10^{-16}$ \\
$\cos(\pi/5)\pket{+}_z + \sin(\pi/5)\mathrm{e}^{i\pi/3}\pket{-}_z$ & $5.55\times10^{-16}$ \\
\bottomrule
\end{tabular}
\end{table}

% ================================================================
\section{The coherence violation indicator and coherence cost}
\label{sec:chi}

\begin{definition}[Coherence violation indicator]
\label{def:chi}
For a mixed triangle $\tau = (e_i, e_j, p_k)$ and a global sign configuration $\sigma$, the coherence violation indicator is:
\begin{equation}
\chi(\tau;\sigma) = s_{\Kfour}(e_i,e_j;\sigma) \cdot s_\times(e_i,p_k;\sigma) \cdot s_\times(e_j,p_k;\sigma) \in \{+1,-1\},
\label{eq:chi_def}
\end{equation}
where $s_{\Kfour}(e_i,e_j;\sigma)$ is the sign of edge $(e_i,e_j)$ in $\Kfour$ under $\sigma$, and $s_\times(e,p;\sigma)$ is the sign of the cross-edge linking $e$ to $p$. We have $\chi(\tau;\sigma) = +1$ if $\tau$ is coherent and $\chi(\tau;\sigma) = -1$ if $\tau$ is violated.
\end{definition}

\begin{remark}[Integer-valued, no probability assumed]
The quantity $\chi(\tau;\sigma) \in \{+1,-1\}$ is a deterministic function of the sign configuration alone. No probability, amplitude, or inner product is assumed in its definition. The violation cost per triangle, $(1-\chi(\tau;\sigma))/2 \in \{0,1\}$, is a non-negative integer: zero for a coherent triangle, one for a violated triangle.
\end{remark}

\begin{definition}[Mixed-triangle coherence cost]
\label{def:emix}
The mixed-triangle coherence cost is~\cite{Born_PDL_2026}:
\begin{equation}
\emix(\sigma;\theta,\psi) = \frac{1}{Z(\theta,\psi)} \sum_{\tau \in \Tplus \cup \Tminus} \alpha_\tau(\theta,\psi)\,\frac{1-\chi(\tau;\sigma)}{2},
\label{eq:emix_def}
\end{equation}
where $\alpha_\tau(\theta,\psi) \geq 0$ are triangle weights and $Z(\theta,\psi) = \sum_\tau \alpha_\tau(\theta,\psi)$ is the normalisation.
\end{definition}

The coherence cost splits into two contributions: the integer counting $(1-\chi)/2 \in \{0,1\}$ is determined by the \PDL\ graph structure alone; the weights $\alpha_\tau$ encode how much each triangle contributes. The derivation of $\alpha_\tau$ from the \PDL\ axioms is open problem~OP1 (Section~\ref{sec:open}).

\begin{definition}[Coherence overlap]
\label{def:ecoh}
The coherence overlap for outcome basin $A$ (outcome $+\theta$) is defined as the coherence cost averaged over the basin:
\begin{equation}
\ecoh(\sigma_+;\theta,\psi) = \emix(\sigma_+;\theta,\psi).
\label{eq:ecoh_def}
\end{equation}
Under the apparatus symmetry conditions of~\cite{Born_PDL_2026}, this evaluates to:
\begin{equation}
\ecoh(\sigma_+;\theta,\psi) = \bigl|\pbra{-\theta}\pket{\psi}\bigr|^2 = \bigl|a\sin(\theta/2) - b\cos(\theta/2)\bigr|^2.
\label{eq:ecoh_id}
\end{equation}
\end{definition}

% ================================================================
\section{The Gleason-type uniqueness theorem (self-contained)}
\label{sec:gleason}

The following theorem is proved in~\cite{Gleason_PDL_2026}. It is reproduced here in full so that the reader can verify the argument of D34 without access to that document.

\begin{axiom}[Operational axioms for spin-$\tfrac{1}{2}$ measurements in \PDL]
\label{ax:five}
Let $P_+(\theta;a,b)$ denote the probability of obtaining outcome $+\theta$ when measuring spin along direction $\theta$ in state $\pket{\psi} = a\pket{+}_z + b\pket{-}_z$, with $|a|^2+|b|^2 = 1$. The following five axioms are imposed:
\begin{enumerate}[label=(\roman*)]
\item \emph{Normalisation}: $P_+(\theta;a,b) + P_-(\theta;a,b) = 1$ for all $\theta, a, b$.
\item \emph{Global phase invariance}: $P_\pm(\theta;\mathrm{e}^{i\phi}a,\mathrm{e}^{i\phi}b) = P_\pm(\theta;a,b)$ for all phases $\phi$.
\item \emph{Rotational covariance}: $P_+(\theta+\theta_0;a,b) = P_+(\theta;a',b')$ where $(a',b') = R_{-\theta_0}(a,b)$ is the state rotated by $-\theta_0$ in $\Hspin$.
\item \emph{Branch symmetry}: exchanging outcome labels $+\leftrightarrow -$ maps $P_+(\theta;a,b)$ to $P_-(\theta+\pi;a,b)$.
\item \emph{Measurement repeatability}: if $\pket{\psi} = \pket{+}_\theta$ then $P_+(\theta;1,0) = 1$.
\end{enumerate}
\end{axiom}

\begin{theorem}[Gleason-type uniqueness for spin-$\tfrac{1}{2}$ in \PDL, \cite{Gleason_PDL_2026}]
\label{thm:gleason}
Under Axioms~\ref{ax:five}(i)--(v), any admissible probability assignment $P_+(\theta;a,b)$ is uniquely determined to be Born's rule:
\begin{equation}
P_+(\theta;a,b) = \bigl|a\cos(\theta/2) + b\sin(\theta/2)\bigr|^2.
\label{eq:born_unique}
\end{equation}
\end{theorem}

The proof of Theorem~\ref{thm:gleason} proceeds as follows (sketch, following~\cite{Gleason_PDL_2026}). Axioms~(i)--(iv) constrain $P_+$ to depend on $(a,b,\theta)$ only through $F(|a\cos(\theta/2)+b\sin(\theta/2)|^2)$ for some continuous function $F:[0,1]\to[0,1]$ satisfying $F(1-x) = 1-F(x)$. Axiom~(v) (repeatability) and the behaviour of $P_+$ on classical mixtures yield Jensen's functional equation for $F$, forcing $F$ to be affine. The symmetry $F(1-x)=1-F(x)$ then forces $F(x) = x$, giving Born's rule uniquely. $\square$

% ================================================================
\section{Born rule from \texorpdfstring{$\ecoh$}{ecoh} and Gleason uniqueness}
\label{sec:born}

\begin{proposition}[Born rule from coherence complement]
\label{prop:born_from_ecoh}
If the coherence overlap satisfies equation~\eqref{eq:ecoh_id}, then:
\begin{equation}
P_+(\theta;\psi) = 1 - \ecoh(\sigma_+;\theta,\psi) = 1 - \bigl|\pbra{-\theta}\pket{\psi}\bigr|^2 = \bigl|\pbra{+\theta}\pket{\psi}\bigr|^2,
\label{eq:born_from_ecoh}
\end{equation}
where the last equality is the normalisation identity $|\pbra{+\theta}\pket{\psi}|^2 + |\pbra{-\theta}\pket{\psi}|^2 = 1$.
\end{proposition}

\begin{proof}
Substituting equation~\eqref{eq:ecoh_id}: $P_+ = 1 - |a\sin(\theta/2) - b\cos(\theta/2)|^2$. Expanding: $|a\sin(\theta/2) - b\cos(\theta/2)|^2 = |a|^2\sin^2(\theta/2) - 2\mathrm{Re}(ab^*)\sin(\theta/2)\cos(\theta/2) + |b|^2\cos^2(\theta/2)$, and $|a\cos(\theta/2) + b\sin(\theta/2)|^2 = |a|^2\cos^2(\theta/2) + 2\mathrm{Re}(ab^*)\cos(\theta/2)\sin(\theta/2) + |b|^2\sin^2(\theta/2)$. Their sum is $(|a|^2+|b|^2)(\sin^2(\theta/2)+\cos^2(\theta/2)) = 1$. Hence $P_+ = |\pbra{+\theta}\pket{\psi}|^2$.
\end{proof}

\begin{proposition}[The five Gleason axioms are satisfied]
\label{prop:axioms}
The assignment $P_+(\theta;\psi) = |\pbra{+\theta}\pket{\psi}|^2$ satisfies all five axioms of Axiom~\ref{ax:five}.
\end{proposition}

\begin{proof}
(i) $|\pbra{+\theta}\pket{\psi}|^2 + |\pbra{-\theta}\pket{\psi}|^2 = \|\pket{\psi}\|^2 = 1$. (ii) $|\pbra{+\theta}(\mathrm{e}^{i\phi}\pket{\psi})|^2 = |\mathrm{e}^{i\phi}|^2|\pbra{+\theta}\pket{\psi}|^2 = |\pbra{+\theta}\pket{\psi}|^2$. (iii) $|\pbra{+(\theta+\theta_0)}\pket{\psi}|^2 = |\pbra{+\theta}R_{-\theta_0}\pket{\psi}|^2 = P_+(\theta; R_{-\theta_0}\pket{\psi})$, verified numerically to $1.1\times10^{-15}$ over 1000 rotations and 20 states (Table~\ref{tab:axioms}). (iv) Exchanging $+\leftrightarrow-$ maps $\pket{+}_\theta\to\pket{-}_\theta$ and $\pket{-}_\theta\to\pket{+}_\theta$, i.e.\ $\theta\to\theta+\pi$; the symmetry holds by orthogonality. (v) $|\pbra{+\theta}\pket{+}_\theta|^2 = 1$.
\end{proof}

\begin{theorem}[Born rule from $(A)\wedge(B)$-admissible coherence]
\label{thm:born}
Under the identification~\eqref{eq:ecoh_id} and the Gleason-type uniqueness Theorem~\ref{thm:gleason}, the probability of obtaining outcome $+\theta$ when measuring the spin of a \PDL\ electron in state $\pket{\psi} = a\pket{+}_z + b\pket{-}_z$ along direction $\theta$ is uniquely:
\begin{equation}
P_+(\theta;a,b) = \bigl|a\cos(\theta/2) + b\sin(\theta/2)\bigr|^2.
\label{eq:born}
\end{equation}
This is Born's rule, and it is the unique assignment compatible with the five Gleason axioms and the \PDL\ coherence structure.
\end{theorem}

\begin{proof}
By Proposition~\ref{prop:born_from_ecoh}, $P_+ = |\pbra{+\theta}\pket{\psi}|^2 = |a\cos(\theta/2)+b\sin(\theta/2)|^2$. By Proposition~\ref{prop:axioms}, this satisfies all five axioms of Axiom~\ref{ax:five}. By Theorem~\ref{thm:gleason}, any assignment satisfying those axioms equals Born's rule. Uniqueness follows.
\end{proof}

\begin{remark}[Epistemic status of Theorem~\ref{thm:born}]
\label{rem:epistemic}
Theorem~\ref{thm:born} is a theorem \emph{under the identification~\eqref{eq:ecoh_id}}, which requires the triangle weights $\alpha_\tau(\theta,\psi)$ to be distributed such that $\emix(\sigma_+;\theta,\psi) = |\pbra{-\theta}\pket{\psi}|^2$. This identification holds under the mild apparatus symmetry conditions of~\cite{Born_PDL_2026} (stated in Section~\ref{sec:born_scheme}). Deriving those conditions directly from the \PDL\ axioms, without invoking the inner-product structure of $\Hspin$, is the residual OP1 of D34 (Section~\ref{sec:open}).
\end{remark}

% ================================================================
\section{Diagnostic: uniform weights lead to Heaviside, not Born}
\label{sec:diagnostic}

\begin{proposition}[Uniform $\alpha_\tau$ leads to Heaviside]
\label{prop:heaviside}
If the triangle weights are uniform, $\alpha_\tau = 1/N_T$ for all $\tau \in \Tplus \cup \Tminus$, then for the eigenstate $\sigma_{\mathrm{sys}} = +1$ the coherence cost is $\emix(+1;\theta) = \sin^2(\theta/2)$ (from equation~\eqref{eq:fractions}), and the Boltzmann basin ratio satisfies:
\begin{equation}
P_A(\theta) = \frac{W_A}{W_A + W_B} = \mathrm{expit}\!\bigl(\lambda N_T\cos\theta\bigr),
\label{eq:sigmoid}
\end{equation}
where $W_A = \exp(-\lambda N_T\sin^2(\theta/2))$ and $W_B = \exp(-\lambda N_T\cos^2(\theta/2))$. In the limit $\lambda N_T \to \infty$, this converges to the Heaviside step function $\Theta(\cos\theta)$, not to $\cos^2(\theta/2)$.
\end{proposition}

\begin{proof}
The log-ratio is $\ln(W_A/W_B) = \lambda N_T(\cos^2(\theta/2) - \sin^2(\theta/2)) = \lambda N_T\cos\theta$. Hence $P_A = \mathrm{expit}(\lambda N_T\cos\theta)$, which converges to $\Theta(\cos\theta)$ as $\lambda N_T \to \infty$.
\end{proof}

\begin{remark}[Structural obstruction, not numerical artefact]
The discrepancy $\Theta(\cos\theta) \neq \cos^2(\theta/2)$ is structural: $\max_\theta|\Theta(\cos\theta) - \cos^2(\theta/2)| = 1$ (attained at $\theta = \pi/2$, see Table~\ref{tab:diagnosis}). This is not a convergence issue; it is a precise structural obstruction showing that uniform weights cannot reproduce the correct geometry of $\Hspin$. Any correct derivation of $\alpha_\tau$ from the \PDL\ axioms must encode the inner-product geometry of $\Hspin$ through the combinatorial structure of $\Kfour \otimes G_{\mathrm{app}}$.
\end{remark}

\begin{table}[ht]
\centering
\caption{Comparison of the two weight models for the Boltzmann basin probability $P_A = W_A/(W_A+W_B)$ in the limit $\lambda \to \infty$. The coherence-overlap model (equation~\eqref{eq:ecoh_id}) converges to Born's rule; the uniform model converges to the Heaviside step function with maximum discrepancy~1.}

\bigskip

\label{tab:diagnosis}
\begin{tabular}{llcr}
\toprule
Model & $\emix(+1;\theta)$ & $\lim_{\lambda\to\infty} P_A(\theta)$ & Max.\ discrepancy from Born \\
\midrule
Uniform $\alpha_\tau = 1/N_T$ & $\sin^2(\theta/2)$ & $\Theta(\cos\theta)$ & $1.00$ \\
Coherence overlap (eq.~\eqref{eq:ecoh_id}) & $|\pbra{-\theta}\pket{\psi}|^2$ & $|\pbra{+\theta}\pket{\psi}|^2$ & $0$ \\
\bottomrule
\end{tabular}
\end{table}

% ================================================================
\section{Numerical checks}
\label{sec:numerics}

All numerical checks were performed in Python\,3 (NumPy\,2.x, SciPy\,1.x). The log-sum-exp technique (\texttt{scipy.special.expit}) was used for numerical stability. Python code is available from the author upon request.

\begin{table}[ht]
\centering
\caption{Numerical verification of the main results of D34. All algebraic identities hold to floating-point precision. The rotational covariance check (Axiom~iii) was performed over 1000 random rotations and 20 random normalised states; the coherence-overlap identity was verified over 1000 states and 20 angles each (20\,000 state--angle pairs total).}
\label{tab:axioms}
\begin{tabular}{llr}
\toprule
Quantity & Expression verified & Max.\ residual \\
\midrule
$(A)\wedge(B) \Leftrightarrow$ stable & 768 cases, 192 stable & 0 counter-examples \\
$\ecoh = |\pbra{-\theta}\pket{\psi}|^2$ & 20\,000 state--angle pairs & $5.55\times10^{-16}$ \\
Axiom~(i): normalisation & 5 states, 500 angles & $5.55\times10^{-16}$ \\
Axiom~(ii): phase invariance & 20 phases, 1 state & $4.44\times10^{-16}$ \\
Axiom~(iii): rotational covariance & 1000 rotations, 20 states & $1.11\times10^{-15}$ \\
Axiom~(iv): branch symmetry & exact algebraic & $0.00\times10^0$ \\
Axiom~(v): repeatability & 3 eigenstates & $0.00\times10^0$ \\
$|\Theta(\cos\theta) - \cos^2(\theta/2)|$ & Heaviside diagnostic & $\leq 1.00$ \\
\bottomrule
\end{tabular}
\end{table}

% ================================================================
\section{Open problems}
\label{sec:open}

\textbf{OP1 (Derivation of $\alpha_\tau$ from \PDL\ axioms).} The central open problem of D34 is the derivation of the triangle weights $\alpha_\tau(\theta,\psi)$ from the \PDL\ axioms alone, without presupposing the inner-product structure of $\Hspin$. Concretely: showing that the combinatorial structure of $\Kfour \otimes G_{\mathrm{app}}$ forces $\sum_{\tau:\,\text{violated in }\sigma_+}\alpha_\tau = |\pbra{-\theta}\pket{\psi}|^2\cdot Z(\theta,\psi)$. The apparatus symmetry conditions of~\cite{Born_PDL_2026} achieve this under external assumptions; deriving those assumptions from the axioms alone remains open.

\textbf{OP2 (Linearity and unitarity from the \PDL\ optimisation functional).} The linear update rule $\psi_{n+1} = a_n\psi_n + b(\psi_{n-1}+\psi_{n+1})$ of D32 is an effective description~\cite{PDL_D32_2026}. Deriving it from the \PDL\ logical optimisation functional would close the last structural gap at the amplitude level.

\textbf{OP3 (Born rule beyond spin-$\tfrac{1}{2}$).} Theorem~\ref{thm:gleason} covers $\dim\Hspin = 2$ only. Extending the Gleason programme to the full \PDL\ lattice of mixed-triangle projectors for $\dim\Hspin \geq 3$ would yield a complete derivation of Born's rule for arbitrary \PDL\ systems.

\textbf{OP4 (Born rule for position measurements).} The amplitudes $\psi_n$ of D32 are position-space amplitudes. Establishing Born's rule for position measurements from the $(A)\wedge(B)$-admissible dynamics of D32 directly, at the same level of rigour as the spin case established here, would connect D34 and D32 at the same generality.

% ================================================================
\section*{Conclusion}

This paper has established the connection between the \PDL\ coherence structure and Born's rule at two epistemic levels, with explicit documentation of what is proved and what remains open.

\emph{At Level~1}, the coherence violation indicator $\chi(\tau;\sigma) \in \{+1,-1\}$ is a purely combinatorial integer, defined from edge signs of $\Kfour \cup G_{\mathrm{app}}$ with no probability assumed. Under identification~\eqref{eq:ecoh_id}, the coherence cost $\ecoh = |\pbra{-\theta}\pket{\psi}|^2$ is the squared overlap of the state with the \emph{orthogonal} outcome direction. Proposition~\ref{prop:born_from_ecoh} gives $P_+ = 1 - \ecoh = |\pbra{+\theta}\pket{\psi}|^2$ as an algebraic consequence of normalisation alone. Proposition~\ref{prop:axioms} verifies the five Gleason axioms for this assignment, and Theorem~\ref{thm:gleason} (Gleason-type uniqueness, stated in full in Section~\ref{sec:gleason}) closes the argument: Born's rule is the unique admissible probability assignment on $\Hspin$ in the \PDL\ setting (Theorem~\ref{thm:born}). This closes open problem~OP5 of D32 at Level~1.

\emph{At Level~2}, Proposition~\ref{prop:heaviside} establishes a structural diagnostic: uniform weights $\alpha_\tau = 1/N_T$ lead to a Heaviside step function with maximum discrepancy~1.00 from Born's rule. This identifies precisely what the weights must encode and defines the target for OP1.

The document is self-contained: Table~\ref{tab:condB} (condition~(B)), Table~\ref{tab:born_check} (Born from coherence complement), and Axiom~\ref{ax:five} + Theorem~\ref{thm:gleason} (Gleason uniqueness) are reproduced in full, so that a reader without access to D29, \cite{Born_PDL_2026}, or \cite{Gleason_PDL_2026} can verify the argument independently.

% ================================================================
\bibliographystyle{plain}
\nocite{*}
\bibliography{References_D34}

\end{document}